% SPDX-License-Identifier: CC-BY-SA-4.0
% Author: Matthieu Perrin
% Part: 
% Section: 
% Sub-section: 
% Frame: 

\begingroup

\begin{frame}[fragile]{Démonstration}

  \vspace{-10mm}
  \begin{lstlisting}[numbers=left, gobble=4]
    private final AtomicInteger range = new AtomicInteger(0);
    private final AtomicReferenceArray<T> items; // Tableau infini

    public void push(T value) {
      int i = range.getAndIncrement();
      items.set(i, value);
    }

    public T pop() {
      for(int i = range.get()-1; i>=0; i--) {
        T value = items.getAndSet(i, null);
        if(value != null) return value;
      }
      return null;
    }
  \end{lstlisting}

  \onBlock<1>[bottom]{Preuve de vivacité}{
    L'algorithme est wait-free
    \begin{description}
    \item[Push :] pas de boucle
    \item[Pop :] au plus \lstinline{range.get()} tours de boucle
    \end{description}
  }%

  \onAlertBlock<2->[bottom]{Exercice : linéarisabilité non-triviale}{%
    Donner deux exécutions :
    \begin{enumerate}
    \item une où le point de linéarisation de push ne peut pas être la ligne 5
    \item une où le point de linéarisation de push ne peut pas être la ligne 6
    \end{enumerate}
  }%

\end{frame}

\endgroup
\endinput
